\section{Conclusion}\label{sec:conclusion}

This study investigated bug resolution dynamics in open-source video games using the OSSGameBench dataset (Version v2). While multi-PR bug issues close substantially later than single-PR issues in observational calendar time (median 35.3 vs.\ 8.0 days, Cliff's $\delta = 0.284$), this raw gap is heavily influenced by time-dependent link accumulation: in multi-PR issues, secondary pull requests appear at a median of 95.3\% of the issue's lifetime. A 30-day fixed-landmark analysis yields a difference of $-0.8$ days that is not statistically distinguishable from zero ($p=0.832$), though this estimate remains inconclusive due to small surviving sample size ($n=33$, 95\% CI $-22.5$ to $29.2$ days). While the landmark design reduces calendar accumulation bias within closed issues, it does not eliminate all selection biases.

Rather than serving as a direct measure of bug difficulty or debugging effort, pull request multiplicity functions primarily as a diagnostic symptom of extended ticket lifecycles. Open-source game maintainers should not treat PR count as a triage proxy for defect complexity, but should instead monitor accumulating PRs as a signal for coordination friction, reviewer churn, or unresolved regressions.
